\documentclass[10pt,a4paper]{amsart}
\usepackage[margin=19mm]{geometry}
\usepackage{amsmath,amssymb,mathtools}
\usepackage[scr=rsfs,cal=euler]{mathalfa}
\usepackage{xfrac}
\usepackage[unicode,hidelinks,bookmarks=false]{hyperref}
\newcommand{\ie}{i.e.\ }
\newcommand{\cf}{cf.\ }
\newcommand{\resp}{respectively\ }
\newcommand{\Subsection}{\subsection}
\providecommand{\nobreakdash}{\nobreak\hbox{-}}
\setlength{\emergencystretch}{2em}
\multlinegap0pt
\usepackage{mathtools}
\usepackage[shortlabels]{enumitem}
\usepackage{xcolor}
\usepackage{amssymb}
\usepackage{csquotes}
\usepackage{cancel}
\usepackage{tikz}
\newtheorem{lemma}{Lemma}
\newcommand{\A}{\mathcal{A}}
\newcommand{\B}{\mathcal{B}}
\newcommand{\C}{\mathcal{C}}
\DeclareMathOperator{\dom}{dom}
\newcommand{\f}{\mathcal{f}}
\title{On the Model Theory of Second-Order Objects}
\author{Tapani Hyttinen, Joni Puljuj\"arvi, Davide Emilio Quadrellaro}
\date{Source: arXiv:2405.03785v2 (CC BY 4.0)}
\begin{document}
\maketitle

\noindent\emph{Source and licence.} On the Model Theory of Second-Order Objects. Version arXiv:2405.03785v2 (CC BY 4.0), distributed by arXiv under the Creative Commons Attribution 4.0 International licence, http://creativecommons.org/licenses/by/4.0/, as recorded in the arXiv licence metadata for this exact version and shown on the abstract page https://arxiv.org/abs/2405.03785v2. Full licence notice: \texttt{licenses/arxiv-2405.03785-CC-BY-4.0.txt}.\par\smallskip

\noindent\textit{Editorial scope.} Counted unit: the article's lemma formulas (source0.tex lines 865--870). The article's complete original proof of it is delivered with it (source0.tex lines 871--980, one \texttt{proof} environment of 110 lines). No mathematics is written, repaired or re-derived: the statement and the proof are the article's own lines, in the article's own notation, and no retained line is substituted or re-flowed.  No further source-local context is needed: every cross-reference of every retained line resolves inside the two delivered spans.  Nothing was omitted from any retained span, so the package carries no omission record.  No retained line contains a citation, so the wrapper carries no bibliography and no bibliography entry.  No retained line contains a figure environment, an image-inclusion command or drawing code, so no image is delivered and the package declares no asset dependency.  Editorial formatting only, in addition: an AMS \texttt{amsart} preamble that restates the article's own declarations and macros used by the retained lines, namely the environments \texttt{lemma} on separate counters, together with the standard packages those lines need.  The retained environments print in this document as Lemma 1 in order of appearance, whereas the article prints the same items with the numbering of its own full document.  The complete native source of the article as distributed in the arXiv e-print for this exact version is bundled unchanged as \texttt{sources/158748/source0.tex}.
\begin{lemma}\label{Partial isomorphisms preserve first-order atomic formulas}
    Let $\f\colon\A\to\B$ be a partial team isomorphism. Then $\f$ preserves atomic first-order formulas, i.e. for any $X\in\dom(\f)$ and atomic first-order formula $\phi(v_0,\dots,v_{n-1})$,
    \[
        \A\models_X\phi \iff \B\models_{\f(X)}\phi.
    \]
\end{lemma}

\begin{proof}
    %One would guess this is relatively easy to see but it proves to be surprisingly tricky due to flatness of first-order logic.

    We begin by showing that for any $\tau$-term $t(v_0,\dots,v_{n-1})$, we have $G_n(t;\A)\in\dom(\f)$ and $\f(G_n(t;\A)) = G_n(t;\B)$, where
    \[
        G_n(t;\C) = \{ (c_0,\dots,c_n)\in\C^{n+1} \mid t^\C(c_0,\dots,c_{n-1}) = c_n \}
    \]
    is the graph of the function $\vec c\mapsto t^\C(\vec c)$. We proceed by induction on $t$.
    \begin{enumerate}
        \item If $t = v_i$, then the function $\vec c\mapsto t^\C(\vec c)$ is the $i^{\text{th}}$ projection. Let $\vec\imath = (0,\dots,n-1,i)$. Then
        \[
            G_n(t;\C) = \{ (c_0,\dots,c_n)\in\C^{n+1} \mid c_i = c_n \} = \Pr_{\vec\imath}(\C^{n+1}).
        \]
        Hence
        \[
            \f(G_n(t;\A)) = \f(\Pr_{\vec\imath}(\A^{n+1})) = \Pr_{\vec\imath}(\f(\A^{n+1})) = \Pr_{\vec\imath}(\B^{n+1}) = G_n(t;\B).
        \]

        \item If $t = c$ for a constant symbol $c\in\tau$, then the function $\vec c\mapsto t^\C(\vec c)$ is the constant map $\vec c \mapsto c$. Then
        \[
            G_n(t;\C) = \{ (c_0,\dots,c_n)\in\C^{n+1} \mid c_n = c^\C \} = \C^n\times\{c^\C\}.
        \]
        Hence
        \begin{align*}
            \f(G_n(t;\A)) &= \f(\A^n\times\{c^\A\}) = \f(\A^n)\times\f(\{c^\A\}) \\
            &= \B^n\times\{c^\B\} = G_n(t;\B).
        \end{align*}

        \item Finally suppose that $t = F(t_0(v_0,\dots,v_{n-1}),\dots,t_{m-1}(v_0,\dots,v_{n-1}))$ for $\tau$-terms $t_i$ and an $m$-ary function symbol $F\in\tau$. By the induction hypothesis, $\f(G_n(t_i;\A)) = G_n(t_i;\B)$ for all $i<m$. First consider the relation
        \[
            X_0(\C) = \{ (t_0^\C(c_0,\dots,c_{n-1}),c_0,\dots,c_{n-1}) \mid c_0,\dots,c_{n-1}\in\C \}.
        \]
        Clearly $X_0(\C) = \Pr_{\vec\imath_0}(G_n(t_0;\C))$ for $\vec\imath_0 = (n,0,\dots,n-1)$, and hence
        \begin{align*}
            \f(X_0(\A)) &= \f(\Pr_{\vec\imath_0}(G_n(t_0;\A))) = \Pr_{\vec\imath_0}(\f(G_n(t_0;\A))) \\
            &= \Pr_{\vec\imath_0}(G_n(t_0;\B)) = X_0(\B).
        \end{align*}
        Then consider the relation
        \[
            X_1(\C) = \{ (t_0^\C(\vec c),t_1^\C(\vec c),c_0,\dots,c_{n-1}) \mid \vec c = (c_0,\dots,c_{n-1})\in\C^n \}.
        \]
        Denoting $\vec\imath_1 = (0,n+1,1,\dots,n)$, we have
        \begin{align*}
            X_1(\C) &= \Pr_{\vec\imath_1}( \{ (t_0^\C(\vec c),c_0,\dots,c_{n-1},t_1^\C(\vec c)) \mid \vec c = (c_0,\dots,c_{n-1})\in\C^n \} ) \\
            &= \Pr_{\vec\imath_1}( X_0(\C) \bowtie_n G_n(t_1;\C) ).
        \end{align*}
        Then
        \begin{align*}
            \f(X_1(\A)) &= \f(\Pr_{\vec\imath_1}( X_0(\A) \bowtie_n G_n(t_1;\A) )) \\
            &= \Pr_{\vec\imath_1}( \f(X_0(\A)) \bowtie_n \f(G_n(t_1;\A)) ) \\
            &= \Pr_{\vec\imath_1}( X_0(\B) \bowtie_n G_n(t_1;\B) ) \\
            &= X_1(\B).
        \end{align*}
        Continuing this way we obtain a description of the relation
        \[
            \hspace{3em} X_{m-1}(\C) = \{ (t_0^\C(\vec c),\dots,t_{m-1}^\C(\vec c),c_0,\dots,c_{n-1}) \mid \vec c = (c_0,\dots,c_{n-1})\in\C^n \}
        \]
        as an element of the closure of $G_n(t_0;\C),\dots,G_n(t_{m-1};\C)$ under projections and join, whence $\f(X_{m-1}(\A)) = X_{m-1}(\B)$. Then, denoting $\vec\jmath_0 = (0,\dots,n-1,n+m)$ and $\vec\jmath_1 = (m,\dots,m+n-1,0,\dots,m-1)$, we have
        \begin{align*}
            \hspace{3em} G_n(t;\C) &= \Pr_{\vec\jmath_0}\left( \{ (\vec c,t_0^\C(\vec c),\dots,t_{m-1}^\C(\vec c),F^\C(t_0^\C(\vec c),\dots,t_{m-1}^\C(\vec c))) \mid \vec c\in\C^n \} \right) \\
            &= \Pr_{\vec\jmath_0}\left(\Pr_{\vec\jmath_1}\big( X_{m-1}(\C)  \big)\bowtie_m F^\C \right),
        \end{align*}
        whence
        \begin{align*}
            \f(G_n(t;\A)) &= \f\left( \Pr_{\vec\jmath_0}\left(\Pr_{\vec\jmath_1}\left( X_{m-1}(\A) \bowtie_m F^\A \right)\right) \right) \\
            &= \Pr_{\vec\jmath_0}\left(\Pr_{\vec\jmath_1}\left( \f(X_{m-1}(\A)) \bowtie_m \f(F^\A) \right)\right) \\
            &= \Pr_{\vec\jmath_0}\left(\Pr_{\vec\jmath_1}\left( X_{m-1}(\B) \bowtie_m F^\B \right)\right) \\
            &= G_n(t;\B).
        \end{align*}
    \end{enumerate}

    Now we show that $\f$ preserves atomic formulas $\phi(v_0,\dots,v_{n-1})$ of first-order logic. Suppose $\phi = R(t_0,\dots,t_{m-1})$ for some $m$-ary relation symbol $R\in\tau$ and $\tau$-terms $t_i(v_0,\dots,v_{n-1})$, $i<m$. Now
    \begin{align*}
        \phi(\C) &\coloneqq \{ (c_0,\dots,c_{n-1})\in\C^n \mid \C\models\phi(c_0,\dots,c_{n-1}) \} \\
        &\,= \{ (c_0,\dots,c_{n-1}) \mid (t_0^\C(c_0,\dots,c_{n-1}),\dots,t_{m-1}^\C(c_0,\dots,c_{n-1}))\in R^\C \} \\
        &\,= \Pr_{(0,\dots,n-1)}( Y^m_n(\C) \bowtie_m R^\C ),
    \end{align*}
    where
    \begin{align*}
        Y^m_n(\C) &\coloneqq \{ (c_0,\dots,c_{n-1},t_0^\C(\vec c),\dots,t_{m-1}^\C(\vec c)) \mid \vec c = (c_0,\dots,c_{n-1})\in\C^n \}.
    \end{align*}
    If we can show that $Y^m_n(\A)\in\dom(\f)$ and $\f(Y^m_n(\A)) = Y^m_n(\B)$, then
    \begin{align*}
        \f(\phi(\A)) &= \f\left( \Pr_{(0,\dots,n-1)}(Y^m_n(\A))\bowtie_m R^\A \right) = \Pr_{(0,\dots,n-1)}( \f(Y^m_n(\A))\bowtie_m\f(R^\A) ) \\
        &= \Pr_{(0,\dots,n-1)}(Y^m_n(\B)\bowtie_m R^\B) \\
        &= \phi(\B).
    \end{align*}
    Then for any $n$-ary $X\in\dom(\f)$,
    \begin{align*}
        \A\models_X\phi &\iff X\subseteq\phi(\A) \iff \f(X) \subseteq \f(\phi(\A)) \\
        &\iff \f(X) \subseteq\phi(\B) \iff \B\models_{\f(X)}\phi,
    \end{align*}
    as desired. So it suffices show that $Y^m_n(\C)$ can be presented in terms of products, projections, intersections and the sets $G(t_i;\C)$. We show this by induction on $m>0$. If $m=1$, then $Y^m_n(\C) = G_n(t_0;\C)$, so we have already proved this. If we have already handled $Y^m_n(\C)$, then
    \begin{align*}
        Y^{m+1}_n(\C) &= \Pr_{\vec\imath}(\{(c_0,\dots,c_{n-1},c_0,\dots,c_{n-1},t_0^\C(\vec c),\dots,t_m^\C(\vec c) \mid \vec c\in\C^n\}) \\
        &= \Pr_{\vec\imath}(\Pr_{\vec\jmath}(Y^m_n(\C)\times G_n(t_m;\C))\cap\Delta^n_\C\times\C^{m+1}),
    \end{align*}
    where $\vec\imath = (0,\dots,n-1,2n,\dots,2n+m)$ (skips the second batch of $c_0,\dots,c_{n-1}$) and $\vec\jmath = (0,\dots,n-1,n+m,\dots,2n+m-1,n,\dots,n+m-1,2n+m)$ (permutes the second batch of $c_0,\dots,c_{n-1}$ to be after the first).

    Then, finally, suppose that $\phi$ is the equation $t_0(v_0,\dots,v_{n-1}) = t_1(v_0,\dots,v_{n-1})$. Now
    \begin{align*}
        \phi(\C) &= \{ (c_0,\dots,c_{n-1})\in\C^n \mid t_0^\C(c_0,\dots,c_{n-1}) = t_1^\C(c_0,\dots,c_{n-1}) \} \\
        &= \Pr_{(0,\dots,n-1)}(G_n(t_0;\C)\cap G_n(t_1;\C)).
    \end{align*}
    Thus $\f(\phi(\A)) = \phi(\B)$. Then we obtain
    \[
        \A\models_X\phi \iff \B\models_{\f(X)}\phi,
    \]
    which concludes the proof. \qedhere
\end{proof}

\end{document}
